\chapter{MirrorOPD 的四种 inner 变体：统一动力学与误差递推}
\label{chap:mirroropd-unified-inner}

本章把 PMD-Part、PMD-Mean、BPO 与 MDPO 放在同一个 outer--inner 框架中。四种方法共享一个冻结的 anchor policy $\pi_k$，区别在于 inner 是否显式表示配分函数、是否使用中心化最优性残差，以及有限样本和有限优化步如何产生误差；PMD-Mean 还会引入总体 target bias。除非特别说明，策略具有完整支持，动作空间有限，teacher 分布记为 $T(\cdot\mid s)$，步长满足 $0<\eta\leq 1$。

\section{统一的 outer 算法}
\label{sec:unified-outer}

第 $k$ 个 outer step 先从 prompt 分布和 anchor policy 采样访问测度：
\[
 x_i\sim\mathcal D,\qquad y_i\sim\pi_k(\cdot\mid x_i),
 \qquad s_{it}=(x_i,y_{i,<t}).
\]
在整个 inner 阶段冻结 $\pi_k$、$T$、访问历史和由它们计算的 target。令
\[
 r_k(s,a)=\log T(a\mid s)-\log\pi_k(a\mid s),
 \qquad
 Z_k(s)=\E_{a\sim\pi_k(\cdot\mid s)}e^{\eta r_k(s,a)}.
\]
精确的 mirror target 是
\begin{equation}
\label{eq:unified-mirror-target}
 \pi_{k+1}^{\star}(a\mid s)
 =\frac{\pi_k(a\mid s)e^{\eta r_k(s,a)}}{Z_k(s)}
 =\frac{\pi_k(a\mid s)^{1-\eta}T(a\mid s)^\eta}{Z_k(s)}.
\end{equation}
outer 更新只在 inner 结束后同步一次：$\pi_{k+1}\leftarrow\pi_{\theta_{k,M_k}}$。如果在 inner 的每一步都把当前 student 写回 anchor，则 KL 参考点移动，算法退化为普通 natural policy gradient，下面关于固定 proximal 子问题的结论不再适用。

\begin{algorithm}
\caption{MirrorOPD 的统一 outer--inner 迭代\\{\small\textit{Semantics: 根据本章统一定义重构；inner 目标由四种变体分别给出}}}
\begin{algorithmic}
\STATE \textbf{Input:} prompt 分布 $\mathcal D$；初始全支持策略 $\pi_0$；步长 $\eta$；outer 次数 $K$；inner 预算 $M_k$
\STATE \textbf{Output:} 策略序列 $(\pi_k)_{k=0}^{K}$
\FOR{$k=0$ \TO $K-1$}
    \STATE 采样 $x_i\sim\mathcal D$、$y_i\sim\pi_k(\cdot\mid x_i)$，构造历史 $s_{it}$
    \STATE 冻结 $\pi_k,T,r_k,Z_k$，设置 $\theta_{k,0}\gets\theta_k$
    \FOR{$m=0$ \TO $M_k-1$}
        \STATE $\theta_{k,m+1}\gets\operatorname{InnerUpdate}(\theta_{k,m};\pi_k,T,r_k)$
    \ENDFOR
    \STATE $\pi_{k+1}(\cdot\mid s)\gets\pi_{\theta_{k,M_k}}(\cdot\mid s)$；$\theta_{k+1}\gets\theta_{k,M_k}$
\ENDFOR
\STATE \textbf{Return:} $(\pi_k)_{k=0}^{K}$
\end{algorithmic}
\end{algorithm}

为简化记号，以下定义 log-ratio
\[
 u_{k,\theta}(s,a)=\log\frac{\pi_\theta(a\mid s)}{\pi_k(a\mid s)},
 \qquad
 c_k(s)=\log Z_k(s).
\]
精确 fixed-point 条件统一为
\begin{equation}
\label{eq:common-fixed-point}
 u_{k,\theta}^{\star}(s,a)=\eta r_k(s,a)-c_k(s),
 \qquad
 \E_{\pi_k}[e^{u_{k,\theta}^{\star}}\mid s]=1.
\end{equation}

\section{PMD-Part：逐动作回归完整 log-ratio}
\label{sec:pmd-part-unified}

PMD-Part 直接回归式 \eqref{eq:common-fixed-point} 的完整 target：
\begin{equation}
\label{eq:pmd-part-objective}
 \mathcal L_{\mathrm{Part},k}(\theta)
 =\frac12\E_{m_k}\!\left[
 u_{k,\theta}(s,a)-\eta r_k(s,a)+c_k(s)
 \right]^2,
\end{equation}
其中 $m_k$ 是 prompt--rollout token 的访问测度。若 $m_k$ 覆盖每个有效 $(s,a)$，零损失当且仅当 $\pi_\theta=\pi_{k+1}^{\star}$。算法实现将 $c_k$ 视为停止梯度的 target；估计 $Z_k$ 时产生的 plug-in 误差属于 outer 误差，而非增加 inner 步数可以消除的优化误差。

\begin{algorithm}
\caption{PMD-Part inner update\\{\small\textit{Semantics: 根据完整 log-ratio target 重构}}}
\begin{algorithmic}
\STATE \textbf{Input:} 冻结 $\pi_k,T$；batch $\mathcal B$；步长 $\alpha$
\STATE \textbf{Output:} 更新后的参数 $\theta^+$
\FOR{$(s,a)\in\mathcal B$}
    \STATE $r\gets\log T(a\mid s)-\log\pi_k(a\mid s)$
    \STATE $c\gets\log\E_{v\sim\pi_k(\cdot\mid s)}[e^{\eta r(s,v)}]$
    \STATE $e\gets u_{k,\theta}(s,a)-\eta r+c$
\ENDFOR
\STATE $\theta^+\gets\theta-\alpha\nabla_\theta\frac{1}{2|\mathcal B|}\sum_{(s,a)\in\mathcal B}e(s,a)^2$
\STATE \textbf{Return:} $\theta^+$
\end{algorithmic}
\end{algorithm}

\section{PMD-Mean：均值近似与 Lambert-$W$ 固定点}
\label{sec:pmd-mean-unified}

这里必须区分两种经常被混用的目标。精确 PMD residual 可以把
$c_k(s)=\log Z_k(s)$ 作为可计算 target；PMD-Mean 则用 anchor 下的 reward
均值代替这个 log-partition。后一替换改变了总体优化问题，即使样本数和 inner
步数都趋于无穷，驻点也不必等于式~\eqref{eq:unified-mirror-target}。

令
\[
 \bar r_k(s)=\E_{\pi_k}[r_k(s,a)\mid s],\qquad
 h_k(s,a)=r_k(s,a)-\bar r_k(s),\qquad \E_{\pi_k}h_k=0.
\]
论文意义下的 PMD-Mean 分布目标是在归一化约束下最小化
\begin{equation}
\label{eq:pmd-mean-objective}
 \mathcal L_{\mathrm{Mean},k}(q)
 =\frac12\E_{a\sim\pi_k(\cdot\mid s)}
 \left[\log\frac{q(a\mid s)}{\pi_k(a\mid s)}-\eta h_k(s,a)\right]^2,
 \qquad \sum_aq(a\mid s)=1.
\end{equation}
精确 target 的平方回归应当在方括号中使用
$\eta r_k-\log Z_k$；PMD-Mean 删除了 action-dependent 的配分函数校正，不能
用“再加一个归一化约束”把二者等同起来。

\begin{theorem}[PMD-Mean 的 Lambert-$W$ 固定点]
\label{thm:four-inner-pmdmean-lambert}
对有限动作、$\pi_k(a\mid s)>0$ 的每个历史，令
$A_k(s)=\E_{\pi_k}e^{\eta h_k}$、$B_k(s)=\E_{\pi_k}e^{2\eta h_k}$。
式~\eqref{eq:pmd-mean-objective} 的唯一分布最优解为
\begin{equation}
\label{eq:four-inner-pmdmean-target}
 q_{\mathrm{Mean},k}(a\mid s)
 =\pi_k(a\mid s)
 \exp\!\left\{\eta h_k(s,a)-W\!\left(\kappa_k(s)e^{\eta h_k(s,a)}\right)\right\},
\end{equation}
其中 $W$ 是主分支，唯一的 $\kappa_k(s)\ge0$ 由归一化决定，并且
\[
 \frac{A_k(A_k-1)}{B_k}\leq\kappa_k\leq\log A_k.
\]
当 $h_k$ 非常数时，$A_k>1$ 且 $\kappa_k>0$。精确 mirror target 与该解的
projective 偏差为
\begin{equation}
\label{eq:pmdmean-target-bias}
 \delta_{\mathrm{Mean},k}(s)
 =\left\|W\!\left(\kappa_ke^{\eta h_k(s,\cdot)}\right)-\log A_k\right\|_{\mathrm{osc}}.
\end{equation}
因此 PMD-Mean 只有在 $h_k$ 为常数，或在小步长极限下该偏差趋于零时，才与
精确 mirror target 重合。
\end{theorem}

\begin{proof}
写 $u(a)=\log(q(a\mid s)/\pi_k(a\mid s))$，引入归一化乘子 $\kappa$，得到
\[
 \mathscr L(u,\kappa)=\frac12\sum_a\pi_k(a\mid s)[u(a)-\eta h_k(a)]^2
 +\kappa\left(\sum_a\pi_k(a\mid s)e^{u(a)}-1\right).
\]
逐 action 求导并除以正的 anchor 概率，得到
\[
 u(a)-\eta h_k(a)+\kappa e^{u(a)}=0.
\]
令 $w_a=\eta h_k(a)-u(a)$，则
$w_ae^{w_a}=\kappa e^{\eta h_k(a)}$，所以
$w_a=W(\kappa e^{\eta h_k(a)})$，代回即得式~\eqref{eq:four-inner-pmdmean-target}。

再定义
\[
 F(\kappa)=\E_{\pi_k}\exp\{\eta h_k-W(\kappa e^{\eta h_k})\}.
\]
函数 $F$ 连续严格递减，$F(0)=A_k\ge1$，且
$F(\kappa)\to0$。故存在唯一根 $F(\kappa_k)=1$；非恒定 $h_k$
时根严格为正。固定此乘子后，拉格朗日函数对 $u$ 的 Hessian 为
\[
 \operatorname{diag}\left(\pi_k(a\mid s)[1+\kappa_ke^{u(a)}]\right),
\]
严格正定，所以该可行驻点是唯一全局最优解。

精确 mirror target 的 log-ratio 是 $\eta h_k-\log A_k$，而
式~\eqref{eq:four-inner-pmdmean-target} 的 log-ratio 是
$\eta h_k-W(\kappa_ke^{\eta h_k})$；两者相减为
$\log A_k-W(\kappa_ke^{\eta h_k})$。去掉 action-independent 常数后，
其振幅正是式~\eqref{eq:pmdmean-target-bias}。
\end{proof}

若代码使用
$-\eta(h_k-\text{stopgrad})\log q$ 的 policy-gradient 线性化损失，
它应标记为 ``centered-PG''，不能直接称为论文的总体 PMD-Mean；单个 sampled
action 也不能无偏恢复 $\bar r_k(s)$。本书实验将这两种实现分开记录，避免把
均值估计误差、Lambert-$W$ bias 与有限 inner 优化误差混在一起。

\begin{algorithm}
\caption{PMD-Mean inner update\\{\small\textit{Semantics: 根据均值近似平方目标重构；不是 exact log-partition 回归}}}
\begin{algorithmic}
\STATE \textbf{Input:} 冻结 $\pi_k,T$；每个历史的 anchor actions $\{a_j\}_{j=1}^{L}$；步长 $\alpha$
\STATE \textbf{Output:} 更新后的参数 $\theta^+$
\FOR{每个历史 $s$}
    \STATE $\bar r\gets L^{-1}\sum_{j=1}^{L}[\log T(a_j\mid s)-\log\pi_k(a_j\mid s)]$
    \FOR{batch 中的 action $a$}
        \STATE $h(s,a)\gets \log T(a\mid s)-\log\pi_k(a\mid s)-\bar r$
        \STATE $u(s,a)\gets \log\pi_\theta(a\mid s)-\log\pi_k(a\mid s)$
        \STATE $e(s,a)\gets u(s,a)-\eta h(s,a)$
    \ENDFOR
\ENDFOR
\STATE $\theta^+\gets\theta-\alpha\nabla_\theta\frac{1}{2|\mathcal B|}\sum_{(s,a)\in\mathcal B}e(s,a)^2$
\STATE \COMMENT{归一化由策略参数化实现；总体极限为 Lambert-$W$ 解}
\STATE \textbf{Return:} $\theta^+$
\end{algorithmic}
\end{algorithm}

\section{BPO：中心化 score 的非线性回归}
\label{sec:bpo-unified}

BPO 同时中心化 reward target 和 student log-ratio。令
\[
 d_{k,\theta}(s,a)=\eta r_k(s,a)-u_{k,\theta}(s,a),
 \qquad
 \bar d_{k,\theta}(s)=\E_{\pi_k}[d_{k,\theta}(s,a)\mid s],
\]
其目标为
\begin{equation}
\label{eq:bpo-objective}
 \mathcal L_{\mathrm{BPO},k}(\theta)
 =\frac{1}{2\eta}\E_{m_k}
 \left[d_{k,\theta}(s,a)-\bar d_{k,\theta}(s)
 \right]^2.
\end{equation}
BPO 的非线性在于 $\bar d_{k,\theta}$ 依赖当前 student。其零点要求 $d_{k,\theta}(s,a)$ 对 action 为常数；结合策略归一化，常数唯一等于 $c_k(s)$，故 fixed point 仍为式 \eqref{eq:common-fixed-point}。有限步时，BPO 的 centered score 会改变 Hessian 和轨迹，虽然精确零点相同。

\begin{algorithm}
\caption{BPO inner update\\{\small\textit{Semantics: 根据当前 student 的中心化 score 重构}}}
\begin{algorithmic}
\STATE \textbf{Input:} 冻结 $\pi_k,T$；anchor actions $\{a_j\}_{j=1}^{L}$；步长 $\alpha$
\STATE \textbf{Output:} 更新后的参数 $\theta^+$
\FOR{每个历史 $s$}
    \FOR{$j=1$ \TO $L$}
        \STATE $d_j\gets\eta[\log T(a_j\mid s)-\log\pi_k(a_j\mid s)]-u_{k,\theta}(s,a_j)$
    \ENDFOR
    \STATE $\bar d\gets L^{-1}\sum_{j=1}^{L}d_j$
    \STATE $\mathcal L_s\gets(2\eta L)^{-1}\sum_{j=1}^{L}(d_j-\bar d)^2$
\ENDFOR
\STATE $\theta^+\gets\theta-\alpha\nabla_\theta\sum_s\mathcal L_s$
\STATE \textbf{Return:} $\theta^+$
\end{algorithmic}
\end{algorithm}

\section{MDPO：KL 正则 surrogate 的直接优化}
\label{sec:mdpo-unified}

MDPO 不回归 log-ratio，而是最大化 KL 正则 surrogate：
\begin{equation}
\label{eq:mdpo-objective}
 \mathcal J_{\mathrm{MDPO},k}(\theta)
 =\E_{s}\left[
 \E_{a\sim\pi_\theta}[r_k(s,a)\mid s]
 -\frac1\eta\KL(\pi_\theta(\cdot\mid s)\Vert\pi_k(\cdot\mid s))
 \right].
\end{equation}
对每个历史，KL 分解给出
\begin{equation}
\label{eq:mdpo-kl-decomposition}
 \mathcal J_{\mathrm{MDPO},k}(\theta\mid s)
 =\frac1\eta\log Z_k(s)
 -\frac1\eta\KL(\pi_\theta(\cdot\mid s)\Vert\pi_{k+1}^{\star}(\cdot\mid s)).
\end{equation}
因此 MDPO 的精确 fixed point 直接是 $\pi_{k+1}^{\star}$。实际实现可以使用完整词表 KL，也可以在 anchor rollout 上使用 action-level importance ratio；后者还需要独立处理 replay history ratio。

\begin{algorithm}
\caption{MDPO inner update\\{\small\textit{Semantics: 根据 KL 正则 surrogate 重构}}}
\begin{algorithmic}
\STATE \textbf{Input:} 冻结 $\pi_k,T$；batch $\mathcal B$；步长 $\alpha$
\STATE \textbf{Output:} 更新后的参数 $\theta^+$
\STATE $\mathcal J\gets0$
\FOR{$s\in\mathcal B$}
    \STATE $\mathcal J\gets\mathcal J+\E_{a\sim\pi_\theta}[r_k(s,a)]-\eta^{-1}\KL(\pi_\theta(\cdot\mid s)\Vert\pi_k(\cdot\mid s))$
\ENDFOR
\STATE $\theta^+\gets\theta+\alpha\nabla_\theta\mathcal J$
\STATE \textbf{Return:} $\theta^+$
\end{algorithmic}
\end{algorithm}

\section{四种 inner 的误差递推与比较定理}
\label{sec:inner-error-comparison}

对每个变体 $v$，先定义其冻结 anchor 子问题的精确分布解
$q_{v,k}$. 对 PMD-Part、BPO 和 MDPO，$q_{v,k}=\pi_{k+1}^{\star}$；对
PMD-Mean，$q_{v,k}=q_{\mathrm{Mean},k}$，由定理
\ref{thm:four-inner-pmdmean-lambert} 给出。令
\[
 e_{v,k}(s)=\left\|\log\frac{\pi_{k+1}(\cdot\mid s)}{q_{v,k}(\cdot\mid s)}\right\|_{\mathrm{osc}},
 \qquad
 \delta_{v,k}(s)=\left\|\log\frac{q_{v,k}(\cdot\mid s)}{\pi_{k+1}^{\star}(\cdot\mid s)}\right\|_{\mathrm{osc}}.
\]
于是 $\delta_{v,k}=0$ 对应 exact mirror 的三种变体，而
$\delta_{\mathrm{Mean},k}$ 正是式~\eqref{eq:pmdmean-target-bias}。

假设冻结目标上的有限 inner 求解满足
\begin{equation}
\label{eq:inner-contraction-assumption}
 e_{v,k}(s)\leq q_{v,k}^{\mathrm{opt}}(s)e_{v,k}^{0}(s)+b_{v,k}(s),
 \qquad 0\leq q_{v,k}^{\mathrm{opt}}(s)<1,
\end{equation}
其中 $e_{v,k}^{0}$ 是 inner 初始策略到 $q_{v,k}$ 的 projective 距离，
$b_{v,k}$ 汇总采样、测度、函数类投影和停止梯度误差。

\begin{theorem}[四种 inner 的统一 outer 误差递推]
\label{thm:four-inner-recursion}
在完整支持、固定 anchor 和式~\eqref{eq:inner-contraction-assumption} 成立时，
四种方法都满足
\begin{equation}
\label{eq:four-inner-recursion}
 \Delta_{k+1}(s)
 \leq (1-\eta)\Delta_k(s)+\delta_{v,k}(s)
 +q_{v,k}^{\mathrm{opt}}(s)e_{v,k}^{0}(s)+b_{v,k}(s),
\end{equation}
其中
$\Delta_k(s)=\|\log(\pi_k(\cdot\mid s)/T(\cdot\mid s))\|_{\mathrm{osc}}$。
若 $e_{v,k}^{0}\leq C_k\Delta_k$，且
$q_{v,k}^{\mathrm{opt}}\leq\bar q$、$C_k\leq\bar C$、
$b_{v,k}\leq\bar b$、$\delta_{v,k}\leq\bar\delta$，并满足
$(1-\eta)+\bar q\bar C<1$，则
\[
 \Delta_K(s)
 \leq[(1-\eta)+\bar q\bar C]^K\Delta_0(s)
 +\frac{\bar\delta+\bar b}{1-(1-\eta)-\bar q\bar C}.
\]
因此 PMD-Mean 的 Lambert-$W$ bias 与有限 inner 误差处在同一个
outer 误差地板中；增加 inner steps 只能减小 $b_{v,k}$，不能消除
$\delta_{\mathrm{Mean},k}$。
\end{theorem}

\begin{proof}
精确 mirror target 满足
\[
 \log\frac{\pi_{k+1}^{\star}(a\mid s)}{T(a\mid s)}
 =(1-\eta)\log\frac{\pi_k(a\mid s)}{T(a\mid s)}-\log Z_k(s).
\]
去掉 action-independent 常数后，得到
$\|\log(\pi_{k+1}^{\star}/T)\|_{\mathrm{osc}}=(1-\eta)\Delta_k$。
另一方面，振幅范数满足三角不等式：
\[
 \Delta_{k+1}
 \leq \left\|\log\frac{\pi_{k+1}^{\star}}T\right\|_{\mathrm{osc}}
 +\left\|\log\frac{q_{v,k}}{\pi_{k+1}^{\star}}\right\|_{\mathrm{osc}}
 +\left\|\log\frac{\pi_{k+1}}{q_{v,k}}\right\|_{\mathrm{osc}}.
\]
代入 $\delta_{v,k}$ 和式~\eqref{eq:inner-contraction-assumption} 即得
式~\eqref{eq:four-inner-recursion}。若再用 $e_{v,k}^{0}\leq C_k\Delta_k$，
便得到标量递推
$\Delta_{k+1}\leq[(1-\eta)+\bar q\bar C]\Delta_k+\bar\delta+\bar b$。
反复代入并求等比数列，得到所述地板界。
\end{proof}

\begin{theorem}[四种 inner 的 fixed-point 分类]
\label{thm:four-inner-fixed-point}
在完整分布参数化、全 action coverage、精确 teacher log-prob 和全局优化条件下：
\begin{enumerate}
    \item PMD-Part 的唯一 fixed point 是 $\pi_{k+1}^{\star}$；
    \item BPO 的中心化最优性残差在归一化后唯一恢复 $\pi_{k+1}^{\star}$；
    \item MDPO 的 KL surrogate 唯一最大化于 $\pi_{k+1}^{\star}$；
    \item 论文意义下的 PMD-Mean 的唯一 fixed point 是
    $q_{\mathrm{Mean},k}$，一般不等于 $\pi_{k+1}^{\star}$，差异由
    $\delta_{\mathrm{Mean},k}$ 精确刻画。
\end{enumerate}
\end{theorem}

\begin{proof}
PMD-Part 的零残差直接给出
$u=\eta r_k-\log Z_k$。BPO 的零残差要求
$d_{k,\theta}(s,a)$ 对 action 为常数；代入
$\E_{\pi_k}e^{u}=1$ 后，该常数被唯一确定为 $\log Z_k$，因此恢复
$u=\eta r_k-\log Z_k$。MDPO 由式~\eqref{eq:mdpo-kl-decomposition} 等于
常数减 $\eta^{-1}\KL(\pi_\theta\Vert\pi_{k+1}^{\star})$，故唯一极值分布是
$\pi_{k+1}^{\star}$。PMD-Mean 的结论则由定理
\ref{thm:four-inner-pmdmean-lambert} 的严格凸拉格朗日问题直接给出；
其 target bias 在 $h_k$ 非常数时通常严格为正。
\end{proof}

\begin{table}[H]
\centering
\caption{四种 inner 的统一比较}
\label{tab:four-inner-comparison}
\begin{tabular}{p{0.17\textwidth}p{0.25\textwidth}p{0.22\textwidth}p{0.24\textwidth}}
\toprule
变体 & inner 目标 & 精确 fixed point & 主要误差源 \\
\midrule
PMD-Part & 完整 log-ratio 平方回归 & $\pi_k^{1-\eta}T^\eta/Z_k$ & $\log Z_k$ plug-in、投影、回归未收敛 \\
PMD-Mean & 均值近似平方回归 & Lambert-$W$ 分布 $q_{\mathrm{Mean},k}$ & 固有 target bias、均值估计、投影与未收敛 \\
BPO & 中心化 student score 的非线性回归 & $\pi_k^{1-\eta}T^\eta/Z_k$ & 当前 student 中心估计、Bellman/teacher 误差 \\
MDPO & reward $-$ reverse KL$/\eta$ & $\pi_k^{1-\eta}T^\eta/Z_k$ & surrogate 优化误差、IS 方差、critic/teacher 误差 \\
\bottomrule
\end{tabular}
\end{table}

最后强调同步纪律：上述定理都要求每个 outer 内固定 $\pi_k$。若在 inner 第 $m$ 步使用 $\pi_{k,m}$ 重新计算 $r$ 或 KL anchor，则目标变成随时间变化的 $\mathcal J_{k,m}$，误差递推需额外加入漂移项
\[
 \zeta_k\;\geq\;\sum_{m=0}^{M_k-1}
 \left\|\log\frac{\pi_{k,m+1}}{\pi_{k,m}}\right\|_{\mathrm{osc}}.
\]
当 $\zeta_k$ 不可求和时，不能直接套用固定 proximal target 的线性收缩结论。
